\section{August 27, 2026}

We continue our discussion of $\sigma$-algebras. Recall that a
$\sigma$-algebra $\mathcal M$ on a set $X$ contains $\varnothing$ and
$X$, is closed under complements, and is closed under countable unions.
If $(X,d)$ is a metric space, then its Borel $\sigma$-algebra
$\mathcal B_X$ is generated by the open subsets of $X$.

\subsection{Product \texorpdfstring{$\sigma$}{sigma}-algebras}

Let $I$ be an index set, let $X_i$ be a nonempty set for each $i\in I$,
and set
\[
  X=\prod_{i\in I}X_i.
\]
For each $i\in I$, let $\mathcal M_i$ be a $\sigma$-algebra on $X_i$,
and let
\[
  \pi_i\colon X\longrightarrow X_i,
  \qquad
  \pi_i((x_j)_{j\in I})=x_i,
\]
be the coordinate projection.

\begin{definition}[Product $\sigma$-algebra]
  The \emph{product $\sigma$-algebra} on $X$ is
  \[
    \bigotimes_{i\in I}\mathcal M_i
    =\mathcal M\left(
      \{\pi_i^{-1}(E):i\in I,\ E\in\mathcal M_i\}
    \right).
  \]
  Thus it is the smallest $\sigma$-algebra that makes every coordinate
  projection measurable.
\end{definition}

The generating sets $\pi_i^{-1}(E)$ are called \emph{cylinder sets}.
They restrict the $i$th coordinate to $E$ and leave all other coordinates
unrestricted.

\begin{lemma}\label{lem:countable-product-rectangles}
  If $I$ is countable, then $\bigotimes_{i\in I}\mathcal M_i$ is
  generated by the measurable rectangles
  \[
    \left\{\prod_{i\in I}E_i:E_i\in\mathcal M_i\right\}.
  \]
\end{lemma}

\begin{proof}
  Since $I$ is countable, every measurable rectangle satisfies
  \[
    \prod_{i\in I}E_i
    =\bigcap_{i\in I}\pi_i^{-1}(E_i)
    \in\bigotimes_{i\in I}\mathcal M_i.
  \]
  Conversely, each cylinder set is a measurable rectangle. Indeed,
  \[
    \pi_j^{-1}(E_j)=\prod_{i\in I}F_i,
    \qquad
    F_i=
    \begin{cases}
      E_j,&i=j,\\
      X_i,&i\neq j.
    \end{cases}
  \]
  The two collections therefore generate the same $\sigma$-algebra.
\end{proof}

The generators of the coordinate $\sigma$-algebras can be used directly
to generate the product.

\begin{proposition}\label{prop:product-generators}
  Suppose
  \[
    \mathcal M_i=\mathcal M(\mathcal G_i),
    \qquad
    \mathcal G_i\subseteq 2^{X_i},
  \]
  for every $i\in I$.
  \begin{enumerate}[label=(\alph*)]
    \item The product $\sigma$-algebra is generated by
      \[
        \mathcal F
        =\{\pi_i^{-1}(G):i\in I,\ G\in\mathcal G_i\}.
      \]
    \item If $I$ is countable and
      $\mathcal G_i^*=\mathcal G_i\cup\{X_i\}$, then the product
      $\sigma$-algebra is generated by
      \[
        \left\{\prod_{i\in I}G_i:G_i\in\mathcal G_i^*\right\}.
      \]
  \end{enumerate}
\end{proposition}

\begin{proof}
  For part (a), every member of $\mathcal F$ belongs to
  $\bigotimes_{i\in I}\mathcal M_i$, so
  \[
    \mathcal M(\mathcal F)
    \subseteq\bigotimes_{i\in I}\mathcal M_i.
  \]
  For the reverse inclusion, fix $i\in I$ and consider
  \[
    \mathcal A_i
    =\{E\subseteq X_i:\pi_i^{-1}(E)\in\mathcal M(\mathcal F)\}.
  \]
  Preimages preserve complements and countable unions, so
  $\mathcal A_i$ is a $\sigma$-algebra. It contains $\mathcal G_i$ and
  hence contains $\mathcal M(\mathcal G_i)=\mathcal M_i$. Therefore
  every cylinder $\pi_i^{-1}(E)$ with $E\in\mathcal M_i$ belongs to
  $\mathcal M(\mathcal F)$, which proves the reverse inclusion.

  For part (b), each cylinder from part (a) is a rectangle whose other
  coordinates are $X_j$. Conversely, since $I$ is countable, every
  rectangle in the displayed collection is the countable intersection
  of the sets $\pi_i^{-1}(G_i)$. When $G_i=X_i$ this set is all of $X$;
  otherwise it is a cylinder from part (a). The conclusion follows.
\end{proof}

\subsection{Products of Borel \texorpdfstring{$\sigma$}{sigma}-algebras}

Let $(X_i,d_i)$, $i=1,\ldots,n$, be metric spaces. On
\[
  X=\prod_{i=1}^n X_i
\]
define the maximum metric
\[
  d(x,y)=\max_{1\leq i\leq n}d_i(x_i,y_i).
\]

\begin{proposition}\label{prop:product-borel-sigma-algebra}
  With the notation above,
  \[
    \bigotimes_{i=1}^n\mathcal B_{X_i}\subseteq\mathcal B_X.
  \]
  If every $(X_i,d_i)$ is separable, then equality holds:
  \[
    \bigotimes_{i=1}^n\mathcal B_{X_i}=\mathcal B_X.
  \]
\end{proposition}

\begin{proof}
  Each coordinate projection $\pi_i\colon X\to X_i$ is continuous. If
  $U_i\subseteq X_i$ is open, then $\pi_i^{-1}(U_i)$ is open in $X$.
  Since $\mathcal B_{X_i}$ is generated by the open subsets of $X_i$,
  Proposition~\ref{prop:product-generators} gives
  \[
    \bigotimes_{i=1}^n\mathcal B_{X_i}\subseteq\mathcal B_X.
  \]

  Now suppose that each $X_i$ is separable, and choose a countable dense
  subset $D_i\subseteq X_i$. The collection
  \[
    \mathcal G_i
    =\{B_{d_i}(p,q):p\in D_i,\ q\in\bb{Q},\ q>0\}
  \]
  is a countable basis for $X_i$ and hence generates
  $\mathcal B_{X_i}$. The finite product
  $D=\prod_{i=1}^nD_i$ is countable and dense in $X$. Moreover,
  \[
    B_d(x,q)=\prod_{i=1}^n B_{d_i}(x_i,q).
  \]
  Thus every ball in $X$ whose center lies in $D$ and whose radius is a
  positive rational belongs to
  $\bigotimes_{i=1}^n\mathcal B_{X_i}$. These balls form a countable
  basis for $X$, so every open subset of $X$ is a countable union of such
  balls. Hence
  \[
    \mathcal B_X\subseteq\bigotimes_{i=1}^n\mathcal B_{X_i},
  \]
  which completes the proof.
\end{proof}

\begin{corollary}
  For every $n\in\bb{N}$,
  \[
    \mathcal B_{\bb{R}^n}
    =\bigotimes_{i=1}^n\mathcal B_{\bb{R}}.
  \]
\end{corollary}

For example, $\mathcal B_{\bb{R}^n}$ is generated by each of the
following collections:
\begin{itemize}
  \item open rectangles $\prod_{i=1}^n(a_i,b_i)$;
  \item closed rectangles $\prod_{i=1}^n[a_i,b_i]$; and
  \item lower orthants $\prod_{i=1}^n(-\infty,a_i]$.
\end{itemize}

\subsection{Measures}

\begin{definition}[Measure]
  Let $\mathcal M$ be a $\sigma$-algebra on a set $X$. A function
  \[
    \mu\colon\mathcal M\longrightarrow[0,\infty]
  \]
  is a \emph{measure} if
  \begin{enumerate}[label=(\roman*)]
    \item $\mu(\varnothing)=0$;
    \item whenever $E_1,E_2,\ldots\in\mathcal M$ are pairwise disjoint,
      \[
        \mu\left(\bigcup_{i=1}^{\infty}E_i\right)
        =\sum_{i=1}^{\infty}\mu(E_i).
      \]
  \end{enumerate}
  The second property is called \emph{countable additivity}. If it is
  required only for finite pairwise disjoint families, then $\mu$ is
  called finitely additive.
\end{definition}

The sets in $\mathcal M$ are called \emph{measurable sets}, and the
triple $(X,\mathcal M,\mu)$ is called a \emph{measure space}. A measure
$\mu$ is
\begin{itemize}
  \item \emph{finite} if $\mu(X)<\infty$;
  \item \emph{$\sigma$-finite} if
    $X=\bigcup_{i=1}^{\infty}E_i$ for some $E_i\in\mathcal M$ with
    $\mu(E_i)<\infty$ for every $i$.
\end{itemize}

A set $N\in\mathcal M$ with $\mu(N)=0$ is called a \emph{null set}. A
statement is said to hold \emph{almost everywhere}, abbreviated a.e., if
it holds outside some null set.

\subsection{Examples of measures}

\begin{example}[Weighted counting measure]
  Let $X$ be nonempty, let $\mathcal M=2^X$, and let
  $f\colon X\to[0,\infty]$. Define
  \[
    \mu(E)=\sum_{x\in E}f(x)
    :=\sup\left\{
      \sum_{x\in F}f(x):F\subseteq E,\ F\text{ finite}
    \right\}.
  \]
  Then $\mu$ is a measure on $2^X$.

  If $f(x)=1$ for every $x\in X$, then $\mu$ is the \emph{counting
  measure},
  \[
    \mu(E)=\#E.
  \]
  If $x_0\in X$ and
  \[
    f(x)=
    \begin{cases}
      1,&x=x_0,\\
      0,&x\neq x_0,
    \end{cases}
  \]
  then
  \[
    \mu(E)=
    \begin{cases}
      1,&x_0\in E,\\
      0,&x_0\notin E.
    \end{cases}
  \]
  This is the \emph{point mass}, or \emph{Dirac measure}, at $x_0$.
\end{example}

\begin{example}[Countable--cocountable measure]
  Let $X$ be uncountable, and let $\mathcal M$ be the
  countable--cocountable $\sigma$-algebra. Then
  \[
    \mu(E)=
    \begin{cases}
      0,&E\text{ is countable},\\
      1,&E^c\text{ is countable}
    \end{cases}
  \]
  defines a measure on $(X,\mathcal M)$.
\end{example}

\begin{example}[Finite additivity is not enough]
  Let $X=\bb{N}$ and $\mathcal M=2^{\bb{N}}$. Define
  \[
    \nu(E)=
    \begin{cases}
      0,&E\text{ is finite},\\
      \infty,&E\text{ is infinite}.
    \end{cases}
  \]
  Then $\nu$ is finitely additive. It is not a measure, because
  \[
    \nu\left(\bigcup_{n=1}^{\infty}\{n\}\right)=\infty
    \qquad\text{but}\qquad
    \sum_{n=1}^{\infty}\nu(\{n\})=0.
  \]
\end{example}

\subsection{Basic properties of measures}

\begin{theorem}\label{thm:basic-measure-properties}
  Let $(X,\mathcal M,\mu)$ be a measure space.
  \begin{enumerate}[label=(\alph*)]
    \item \emph{Monotonicity:} if $E,F\in\mathcal M$ and $E\subseteq F$,
      then
      \[
        \mu(E)\leq\mu(F).
      \]
    \item \emph{Countable subadditivity:} if
      $E_1,E_2,\ldots\in\mathcal M$, then
      \[
        \mu\left(\bigcup_{i=1}^{\infty}E_i\right)
        \leq\sum_{i=1}^{\infty}\mu(E_i).
      \]
    \item \emph{Continuity from below:} if
      \[
        E_1\subseteq E_2\subseteq\cdots,
        \qquad E_i\in\mathcal M,
      \]
      then
      \[
        \mu\left(\bigcup_{i=1}^{\infty}E_i\right)
        =\lim_{i\to\infty}\mu(E_i).
      \]
    \item \emph{Continuity from above:} if
      \[
        E_1\supseteq E_2\supseteq\cdots,
        \qquad E_i\in\mathcal M,
        \qquad \mu(E_1)<\infty,
      \]
      then
      \[
        \mu\left(\bigcap_{i=1}^{\infty}E_i\right)
        =\lim_{i\to\infty}\mu(E_i).
      \]
  \end{enumerate}
\end{theorem}

\begin{proof}
  For part (a), write $F$ as the disjoint union
  \[
    F=E\cup(F\setminus E).
  \]
  Countable additivity gives
  \[
    \mu(F)=\mu(E)+\mu(F\setminus E)\geq\mu(E).
  \]

  For part (b), define
  \[
    F_1=E_1,
    \qquad
    F_k=E_k\setminus\bigcup_{j=1}^{k-1}E_j
    \quad(k\geq2).
  \]
  The sets $F_k$ are pairwise disjoint, $F_k\subseteq E_k$, and
  \[
    \bigcup_{k=1}^{\infty}F_k
    =\bigcup_{k=1}^{\infty}E_k.
  \]
  Therefore countable additivity and part (a) imply
  \[
    \mu\left(\bigcup_{k=1}^{\infty}E_k\right)
    =\sum_{k=1}^{\infty}\mu(F_k)
    \leq\sum_{k=1}^{\infty}\mu(E_k).
  \]

  For part (c), set
  \[
    F_1=E_1,
    \qquad
    F_k=E_k\setminus E_{k-1}
    \quad(k\geq2).
  \]
  The sets $F_k$ are pairwise disjoint, and for each $N$,
  \[
    E_N=\bigcup_{k=1}^N F_k,
    \qquad
    \bigcup_{i=1}^{\infty}E_i
    =\bigcup_{k=1}^{\infty}F_k.
  \]
  It follows that
  \begin{align*}
    \mu\left(\bigcup_{i=1}^{\infty}E_i\right)
    &=\sum_{k=1}^{\infty}\mu(F_k)\\
    &=\lim_{N\to\infty}\sum_{k=1}^N\mu(F_k)\\
    &=\lim_{N\to\infty}\mu(E_N).
  \end{align*}

  For part (d), define
  \[
    F_i=E_1\setminus E_i.
  \]
  Since $E_i$ decreases, $F_i$ increases, and
  \[
    \bigcup_{i=1}^{\infty}F_i
    =E_1\setminus\bigcap_{i=1}^{\infty}E_i.
  \]
  By part (c),
  \[
    \mu\left(E_1\setminus\bigcap_{i=1}^{\infty}E_i\right)
    =\lim_{i\to\infty}\mu(E_1\setminus E_i).
  \]
  Because $E_i\subseteq E_1$ and $\mu(E_1)<\infty$, finite additivity
  gives
  \[
    \mu(E_1\setminus E_i)=\mu(E_1)-\mu(E_i)
  \]
  and
  \[
    \mu\left(E_1\setminus\bigcap_{i=1}^{\infty}E_i\right)
    =\mu(E_1)-\mu\left(\bigcap_{i=1}^{\infty}E_i\right).
  \]
  Substituting these identities into the preceding limit and cancelling
  the finite quantity $\mu(E_1)$ yields
  \[
    \mu\left(\bigcap_{i=1}^{\infty}E_i\right)
    =\lim_{i\to\infty}\mu(E_i).
  \]
\end{proof}

\begin{remark}[Necessity of the finite-measure hypothesis]
  The assumption $\mu(E_1)<\infty$ in
  Theorem~\ref{thm:basic-measure-properties}(d) cannot be omitted. Let
  \[
    X=\bb{N},
    \qquad
    \mathcal M=2^{\bb{N}},
  \]
  let $\mu$ be counting measure, and define
  \[
    E_i=\{i,i+1,i+2,\ldots\}.
  \]
  Then
  \[
    E_1\supseteq E_2\supseteq E_3\supseteq\cdots,
    \qquad
    \bigcap_{i=1}^{\infty}E_i=\varnothing.
  \]
  Indeed, for any $n\in\bb{N}$, we have $n\notin E_i$ whenever $i>n$.
  Consequently,
  \[
    \mu\left(\bigcap_{i=1}^{\infty}E_i\right)
    =\mu(\varnothing)=0.
  \]
  On the other hand, every $E_i$ is infinite, so
  \[
    \mu(E_i)=\infty
    \quad\text{for every }i,
    \qquad
    \lim_{i\to\infty}\mu(E_i)=\infty.
  \]
  Thus
  \[
    \mu\left(\bigcap_{i=1}^{\infty}E_i\right)
    =0\neq\infty
    =\lim_{i\to\infty}\mu(E_i).
  \]
  This does not contradict part (d), because
  $\mu(E_1)=\mu(\bb{N})=\infty$.
\end{remark}
